\documentclass[12pt]{article}
\usepackage[T1]{fontenc}
\usepackage[utf8]{inputenc}
\usepackage{lmodern}
\usepackage{textcomp}
\usepackage{enumitem}
\usepackage{geometry}
\usepackage{graphicx}
\usepackage{amsmath, amssymb}
\geometry{margin=1in}
\begin{document}
\begin{center}
History - K-12 Level Assessment 18 \
Total Marks: 40 \
Answer all questions.
\end{center}

\section*{Multiple Choice (10 marks)}
\begin{enumerate}[label=\arabic*.]
\item In the author's two-hour movie representing the history of the universe, when do the earliest members of the human family appear?
\begin{enumerate}[label=(\alph*)]
    \item at the beginning of the film
    \item about half way through the film
    \item 90 minutes into the film
    \item during the final 3 seconds of the film
\end{enumerate}
\item The ridge of bone that runs along the top of the skull in Australopithecus robustus is called a:
\begin{enumerate}[label=(\alph*)]
    \item diastema.
    \item cranium.
    \item maxilla.
    \item sagittal crest.
\end{enumerate}
\item The royal graves of the Shang Dynasty consisted of enormous cruciform-shaped tombs, where the deceased kings were buried with:
\begin{enumerate}[label=(\alph*)]
    \item jade, bronze, and ceramic artifacts.
    \item chariots and horses.
    \item the headless bodies of humans who had been sacrificed.
    \item all of the above.
\end{enumerate}
\item The idea that cultures change through time in relation to their environments, along many different paths, is known as:
\begin{enumerate}[label=(\alph*)]
    \item unilineal evolution.
    \item environmental uniformitarianism.
    \item multicultural adaptation.
    \item multilineal evolution.
\end{enumerate}
\item Machu Picchu was built in the mountains as a:
\begin{enumerate}[label=(\alph*)]
    \item defensive stronghold.
    \item shrine to the Inca rain god.
    \item relaxing getaway for elites.
    \item astronomical observatory.
\end{enumerate}
\end{enumerate}

\section*{True / False (10 marks)}
\textit{Answer True or False}
\begin{enumerate}[label=\arabic*.]
\setcounter{enumi}{5}
\item True or False: A person who believes the Earth's features result from natural disasters is known as a catastrophist.
\item True or False: The Denali Complex is a lithic technology used in the Arctic featuring micro-blades, bifacial knives, and burins.
\item True or False: Mayan city-states were organized similarly to Egyptian communities, sharing common governance structures.
\item True or False: The mean cranial capacities for premodern Homo sapiens, Neandertals, and modern Homo sapiens are 1261 cc, 1480 cc, and 1450 cc respectively.
\item True or False: The first hominids are accurately described as bipedal apes, capable of walking on two legs.
\end{enumerate}

\section*{Long Form Response (20 marks)}
\textit{Answer each question with a clear and complete explanation.}
\begin{enumerate}[label=\arabic*.]
\setcounter{enumi}{10}
\item Discuss how burials with exotic grave goods reflect social stratification and inequality in ancient societies, considering the implications of labor intensity on social status.
\item Discuss how archaeologists characterize the societies that constructed the monumental structures at Göbekli Tepe, Watson Brake, and Poverty Point, focusing on their lifestyle and resource management.
\end{enumerate}

\vfill
\noindent\textit{End of Paper}
\end{document}